% ============================================================================
% 论文 RTX 4090 D 全量替换包（自动生成，2026-08）
% 应用方式：按节替换 manuscript/sections/*.tex 对应段落/表格。
% 口径：全稿统一默认 HOPE 参数（theta=0.5, eps=0.02）；已删除优化参数表
%       （tab:multiseed-optimized），避免混合口径。
% 数据源：result/ablation_suite_4090_rerun | ablation_4090_multiseed |
%         scene_out_4090 | ablation_reflection_only_4090（runtime_manifest 实测）
% ============================================================================

% ----------------------------------------------------------------------------
% 1) experimental_setup.tex:79 硬件与运行元数据（替换整段 "The refreshed
%    canonical manuscript evidence was rerun ..."）
% ----------------------------------------------------------------------------
\paragraph{替换文本：}
The refreshed canonical manuscript evidence was rerun with the local backend
model \texttt{deepseek-r1-0528-qwen3-8b} inside a dedicated Python environment.
The recorded runtime manifests report Python 3.12.7, PyTorch 2.5.1, an
OpenAI-compatible local API endpoint, a 13th Gen Intel(R) Core(TM) i7-13700KF
CPU, and an NVIDIA GeForce RTX 4090 D GPU (24\,GB). The Python-side torch build
is CPU-only; GPU inference is served by the local OpenAI-compatible endpoint.
Newly generated \texttt{full\_result.json} files emit this information through a
top-level \texttt{runtime\_manifest} ...

% ----------------------------------------------------------------------------
% 2) discussion.tex:22 硬件句（替换 "...i7-13700F CPU with an NVIDIA GeForce
%    RTX 3060 GPU..." 为下句）
% ----------------------------------------------------------------------------
\paragraph{替换文本：}
... on a 13th Gen Intel(R) Core(TM) i7-13700KF CPU with an NVIDIA GeForce RTX
4090 D GPU (24\,GB), using an OpenAI-compatible local endpoint.

% ----------------------------------------------------------------------------
% 3) Table 4（tab:ablation-main）替换（4090 默认参数，含 Reflection 行）
%    caption 改为：All groups use default HOPE parameters ($\theta=0.5$,
%    $\epsilon=0.02$). The reflection-only row uses the matched follow-up rerun.
% ----------------------------------------------------------------------------
\begin{table}[ht]
\centering
\caption{Main ablation results on the sample coastal joint-assault scenario
(RTX 4090 D). All groups use 10 planning iterations, 50 Monte Carlo runs,
seed 7, writeback disabled, and default HOPE parameters
($\theta=0.5$, $\epsilon=0.02$).}
\label{tab:ablation-main}
\small
\setlength{\tabcolsep}{3.2pt}
\begin{tabular}{lccccc}
\toprule
Setting & MS & OE & LER & Time (h) & Cmd.\ resilience \\
\midrule

Pure LLM & 62.75 & 73.84 & 1.71 & 11.41 & 60.96 \\
LLM + Memento & 61.25 & 70.90 & 1.52 & 13.48 & 60.96 \\
LLM + Hope & 64.45 & 77.29 & 1.97 & 10.92 & 63.59 \\
LLM + Reflection & 59.29 & 67.49 & 1.17 & 11.56 & 60.96 \\
Full & 65.76 & 78.75 & 2.19 & 14.18 & 76.08 \\
\bottomrule
\end{tabular}
\end{table}

% ----------------------------------------------------------------------------
% 4) Table 5（tab:ablation-stage）替换（列序：Stage & Pure & Memento & Hope &
%    Full & Reflection）
% ----------------------------------------------------------------------------
\begin{table}[ht]
\centering
\caption{Five-stage kill-chain scores for the ablation study (RTX 4090 D).}
\label{tab:ablation-stage}
\footnotesize
\setlength{\tabcolsep}{4pt}
\begin{tabular}{lcccccc}
\toprule
Stage & Pure LLM & \shortstack{LLM +\\Memento} & \shortstack{LLM +\\Hope} & Full & \shortstack{LLM +\\Reflection} \\
\midrule

detect & 62.24 & 61.00 & 64.94 & 66.65 & 61.59 \\
disrupt & 57.31 & 61.26 & 57.78 & 65.69 & 60.25 \\
breach & 55.91 & 55.83 & 58.57 & 60.88 & 52.41 \\
control & 64.43 & 60.88 & 66.50 & 67.97 & 54.21 \\
sustain & 68.40 & 64.03 & 70.09 & 66.57 & 55.22 \\
\bottomrule
\end{tabular}
\setlength{\tabcolsep}{6pt}
\end{table}

% ----------------------------------------------------------------------------
% 5) Table 12（tab:multiseed）替换（4090 5-seed 配对检验）
%    叙述：mean MS gain +3.81 (SD=2.27), t(4)=3.75, p=0.02, d=1.68;
%          OE gain +4.85 (t=5.87, p<0.01); LER gain +0.486 (t=6.60, p<0.01)
% ----------------------------------------------------------------------------
\begin{table}[ht]
\centering
\caption{Multi-seed replication (RTX 4090 D): Pure LLM vs.\ Full pipeline
across five seeds. All runs use 10 iterations, 50 Monte Carlo runs, default
HOPE parameters, and writeback disabled.}
\label{tab:multiseed}
\small
\begin{tabular}{lccc}
\toprule
Seed & Pure LLM & Full & $\Delta$MS \\
\midrule

7 & 63.33 & 66.92 & +3.59 \\
42 & 63.92 & 65.96 & +2.04 \\
123 & 63.97 & 71.68 & +7.71 \\
999 & 62.73 & 65.99 & +3.26 \\
2024 & 65.64 & 68.08 & +2.44 \\
\bottomrule
\end{tabular}
\end{table}

% ----------------------------------------------------------------------------
% 6) Table 6/7（tab:generalization-summary / tab:generalization-scenes）
% ----------------------------------------------------------------------------

% ===== Table 6 (tab:generalization-summary) — 4090 scene-out suite =====
% 场景数 5 | Full 胜 Pure 4/5 | Full 不低于全部变体 2/5
\begin{table}[ht]
\centering
\caption{Cross-scenario transfer summary (4090, 20-iteration scene-out suite). All settings use default HOPE parameters ($\theta=0.5$, $\epsilon=0.02$).}
\label{tab:generalization-summary}
\small
\begin{tabular}{lcc}
\toprule
Setting & Avg. MS & Avg. OE \\
\midrule
Pure LLM & 65.01 & 72.49 \\
LLM + Memento & 65.46 & 72.93 \\
LLM + Hope & 66.40 & 78.01 \\
Full & 66.48 & 77.92 \\
\bottomrule
\end{tabular}
\end{table}

% ===== Table 7 (tab:generalization-scenes) — 4090 per-scene =====
\begin{table}[ht]
\centering
\caption{Per-scene mission success (4090, 20-iteration scene-out suite).}
\label{tab:generalization-scenes}
\small
\setlength{\tabcolsep}{4pt}
\begin{tabular}{lcccccc}
\toprule
Scenario & Family & Pure & Memento & Hope & Full & $\Delta$MS \\
\midrule
Coastal joint assault & assault & 58.98 & 58.55 & 60.84 & 60.44 & +0.01 \\
Island resupply corridor & sustain & 64.87 & 64.00 & 65.28 & 65.93 & +0.01 \\
Mountain corridor recon & recon & 68.82 & 69.85 & 71.60 & 71.71 & +0.03 \\
River crossing breakthrough & assault & 64.84 & 67.48 & 69.19 & 67.84 & +0.03 \\
Urban hub defense & defense & 67.52 & 67.40 & 65.11 & 66.50 & -0.01 \\
\bottomrule
\end{tabular}
\end{table}

% ===== 叙述要点（results.tex 段落替换参考） =====
% - Full 胜 Pure：4/5 场景；平均 MS Pure=65.01 Full=66.48（Δ+1.48 点）
% - Full 不低于全部变体：2/5 场景
% - 若有场景非 Full 最优，列出（如 river 场景 Hope 领先），如实报告



% ----------------------------------------------------------------------------
% 7) 摘要（abstract.tex）重写要点
% ----------------------------------------------------------------------------
% - 单场景：Full 相对 Pure：MS 62.75 -> 65.76（Δ+3.01 点），OE 73.84 -> 78.75
%   （Δ+4.91 点），LER 1.71 -> 2.19（4090 默认参数）。
% - 多种子：5 seeds 配对 ΔMS = +3.81（SD=2.27），t(4)=3.75，p=0.02，
%   Cohen's d=1.68；ΔOE t(4)=5.87 p<0.01；ΔLER t(4)=6.60 p<0.01。
% - 泛化：5 场景中 Full 胜 Pure 4/5，平均 MS 65.01 -> 66.48（Δ+1.48 点），
%   平均 OE 72.49 -> 77.92（Δ+5.43 点）；coastal/river 场景 Hope 领先，
%   urban 场景 Pure 领先（seed 7 单次运行，如实报告）。
% - 删除所有 "optimized HOPE parameters (theta=0.9, eps=0.005)" 表述。

% ----------------------------------------------------------------------------
% 8) results.tex 叙述改写要点
% ----------------------------------------------------------------------------
% - Table 4 段：Full 相对 Pure 提升改为 +3.01 点 MS、+4.91 点 OE；删除
%   "The Full setting uses optimized HOPE SDE parameters" 表述。
% - Table 5 段：如实改写为 "Full 在 detect/disrupt/breach/control 四个阶段
%   领先，sustain 阶段 Hope 保持单项优势（70.09 vs 66.57）——该现象与
%   breach/sustain 竞争 protection 维度的结构性张力一致（见 4.5 节）"。
% - 泛化段：Full 胜 Pure 由 "all five scenarios" 改为 "4/5 scenarios
%   （urban-defense 除外）"；平均增益 +1.48 点；说明单 seed 单次运行方差。
% - 删除 tab:multiseed-optimized 表及 "optimized parameters" 相关段落。
% - Figure 数据（ablation_stage_profile.png 等）需按新表重新生成
%   （scripts/generate_paper_figures.py 数据源切换为 4090 套件）。

